% META: {"id": "TSPE-GEOESPACE-CR-011", "chapitre": "TSPE-GEOMETRIE-ESPACE", "type_objet": "cours", "capacites_codes": ["C3", "C6"], "sources_inspiration": [], "mode_creation": "ex_nihilo", "status": "approved"}

\section{Positions relatives de droites et de plans}

% ============================================================
% STRATE 1 — Cas de figure
% ============================================================

\propriete[Deux droites de l'espace]{
Deux droites de l'espace peuvent être : \textbf{coplanaires} (secantes, ou strictement paralleles, ou confondues) ou \textbf{non coplanaires} (on dit aussi non secantes et non paralleles). Contrairement au plan, deux droites de l'espace qui ne sont pas paralleles ne sont pas necessairement secantes.
}

\propriete[Droite et plan]{
Une droite $\mathcal{D}$ et un plan $\mathcal{P}$ de l'espace peuvent être : \textbf{secants} en un unique point, \textbf{strictement paralleles} (aucun point commun), ou $\mathcal{D}$ est \textbf{incluse} dans $\mathcal{P}$.
}

\propriete[Deux plans]{
Deux plans de l'espace sont soit \textbf{secants} (leur intersection est une droite), soit \textbf{strictement paralleles} (aucun point commun), soit \textbf{confondus}.
}

\definition[1]{
Des points ou des vecteurs sont \textbf{coplanaires} s'ils appartiennent a un même plan. Trois vecteurs $\vec{u},\vec{v},\vec{w}$ sont coplanaires si et seulement si il existe des réels $a,b$ tels que $\vec{w}=a\vec{u}+b\vec{v}$ (en supposant $\vec{u},\vec{v}$ non colineaires).
}

\exemple{
Dans un tetraedre $ABCD$, les droites $(AB)$ et $(CD)$ sont-elles coplanaires ?

Si elles etaient coplanaires, les quatre points $A,B,C,D$ seraient coplanaires, ce qui contredit la définition d'un tetraedre (ses quatre sommets ne sont jamais coplanaires). Les droites $(AB)$ et $(CD)$ sont donc non coplanaires.
}

% ============================================================
% STRATE 2 — Erreurs frequentes
% ============================================================

\erreurFrequente{
\textbf{Croire que deux droites non paralleles sont toujours secantes.} C'est vrai dans le plan mais faux dans l'espace : deux droites non paralleles et non secantes existent (droites non coplanaires).
}

% BEGIN-VERIFY
% from sympy import Matrix, symbols, solve, simplify
%
% t, s = symbols('t s', real=True)
%
% # Le tetraedre de l'exemple : quatre points non coplanaires.
% A, B, C, D = (0, 0, 0), (1, 0, 0), (0, 1, 0), (0, 0, 1)
% def vecteur(U, V): return tuple(V[n] - U[n] for n in range(3))
% assert Matrix([vecteur(A, B), vecteur(A, C), vecteur(A, D)]).det() != 0
%
% # (AB) et (CD) ne sont donc pas coplanaires : le systeme d'intersection n'a
% # pas de solution, et leurs directions ne sont pas colineaires.
% AB, CD = vecteur(A, B), vecteur(C, D)
% assert Matrix([AB, CD]).rank() == 2                     # non paralleles
% point_AB = tuple(A[n] + t*AB[n] for n in range(3))
% point_CD = tuple(C[n] + s*CD[n] for n in range(3))
% assert solve([point_AB[n] - point_CD[n] for n in range(3)], [t, s]) == []
%
% # Les trois positions d'une droite et d'un plan, sur le plan z = 0.
% directions = {
%     "secante": ((0, 0, 1), (0, 0, 5)),        # traverse
%     "parallele": ((1, 0, 0), (0, 0, 5)),      # au-dessus, jamais atteint
%     "incluse": ((1, 0, 0), (0, 0, 0)),        # dans le plan
% }
% attendus = {"secante": 1, "parallele": 0, "incluse": None}
% for nom, (direction, origine) in directions.items():
%     solutions = solve(origine[2] + t*direction[2], t)
%     if attendus[nom] == 1:
%         assert len(solutions) == 1
%     elif attendus[nom] == 0:
%         assert solutions == []
%     else:
%         assert solutions == []      # 0 = 0 : sympy ne renvoie pas de racine
%         assert origine[2] == 0 and direction[2] == 0    # tout le monde convient
%
% # Trois vecteurs coplanaires : w s'ecrit comme combinaison de u et v.
% u, v = (1, 0, 0), (0, 1, 0)
% w = (3, -2, 0)
% assert Matrix([u, v, w]).det() == 0
% a, b = symbols('a b', real=True)
% assert solve([a*u[n] + b*v[n] - w[n] for n in range(3)], [a, b]) == {a: 3, b: -2}
% # Et un vecteur hors du plan ne s'y ecrit pas.
% hors = (0, 0, 1)
% assert Matrix([u, v, hors]).det() != 0
% assert solve([a*u[n] + b*v[n] - hors[n] for n in range(3)], [a, b]) == []
% END-VERIFY
